\chapter{RESULTADOS}\label{cap4Resultados}

Este capítulo apresenta os resultados obtidos na validação do ERP Onmarmoraria: testes funcionais do painel e da API em ambiente local, testes manuais em dispositivos móveis reais (iPhone e Android), medições de desempenho sob carga, indicadores objetivos de qualidade de código e a verificação da etapa de integração contínua (\textit{build} e publicação de imagem). A discussão desses resultados, o alcance dos objetivos e as limitações observadas são apresentados no Capítulo~\ref{cap5Discussao}.

\section{Links do Projeto}

Os repositórios e componentes principais do projeto estão disponíveis em:

\begin{itemize}
  \item Frontend (painel): \url{https://github.com/marcelobracet/erp-dashboard-v1}
  \item Backend (API): \url{https://github.com/marcelobracet/erp-api-tcc}
  \item Landing page institucional: \url{https://github.com/marcelobracet/lp-erp-marmoraria}
\end{itemize}

Como discutido no Capítulo~\ref{cap5Discussao}, optou-se por não manter uma instância pública em produção durante este trabalho, de modo que os links acima referem-se exclusivamente ao código-fonte versionado, garantindo que a documentação permaneça acessível independentemente do ciclo de vida de qualquer infraestrutura de nuvem.

\section{Ambiente de Testes}

Todos os testes descritos neste capítulo foram conduzidos em ambiente local (macOS, Docker Compose para API e PostgreSQL, e o painel executado via contêiner Docker ou \texttt{next dev}), reproduzindo a mesma topologia de contêineres prevista para a arquitetura em nuvem descrita nos Capítulos~\ref{cap2FundamentosTecnicos} e~\ref{cap3MatMet}. Essa escolha, motivada pelo caráter de MVP acadêmico e pela necessidade de manter o material deste trabalho sempre acessível durante o período de avaliação e defesa, sem depender da disponibilidade contínua de uma máquina virtual paga, é detalhada no Capítulo~\ref{cap5Discussao}.

A camada de infraestrutura em nuvem prevista para o projeto (máquina virtual \texttt{e2-micro} na Google Cloud Platform ou VPS equivalente, \textit{proxy} reverso Nginx com TLS via Let's Encrypt e orquestração via Kubernetes para cenários de maior escala) foi implementada e versionada no repositório (\texttt{docker-compose.prod.yml}, manifestos de Kubernetes e script de \textit{bootstrap} da VM, Seção~\ref{sec:workflow-deploy}), mas não foi exercitada contra uma instância pública nesta fase do trabalho. Essa decisão implica uma limitação relevante, discutida na Seção~\ref{sec:limitacoes} do Capítulo~\ref{cap5Discussao}: riscos de paridade entre ambiente local e ambiente de nuvem real (por exemplo, escopos de permissão de IAM, variáveis de ambiente específicas de produção e latência de rede) não puderam ser exercitados de ponta a ponta neste trabalho.

\section{Testes Funcionais}
\label{sec:testes-funcionais}

\subsection{Autenticação, \textit{Provision} e Isolamento por \textit{Tenant}}

O fluxo de \textit{login} com e-mail e senha, emissão de JWT e renovação silenciosa via \textit{refresh token} foi validado em ambiente local. O \textit{provision} de uma nova empresa (criação de \textit{tenant}, usuário administrador e configurações iniciais em uma única operação) foi exercitado repetidas vezes durante os testes, inclusive diretamente contra a API local via linha de comando, confirmando a emissão correta do \textit{token} e das permissões associadas ao papel \textit{admin}. O isolamento por \texttt{tenant\_id} foi verificado criando registros em diferentes \textit{tenants} e confirmando que cada um enxerga apenas os seus próprios dados.

\subsection{Dashboard e Módulos de Negócio}

O painel exibiu corretamente os indicadores de orçamentos, conversão, faturamento estimado e ticket médio, calculados no navegador a partir dos dados retornados pela API (os valores atuais utilizam dados de demonstração, conforme indicado na própria interface). Os módulos de clientes, produtos e orçamentos foram exercitados com operações de criação, leitura, atualização e exclusão autenticadas.

\subsection{Página Institucional e Lista de Espera}

A página institucional (\textit{landing page}) foi validada quanto à navegação e ao envio do formulário de lista de espera. Com a API em execução, o \texttt{POST /\allowbreak{}api/\allowbreak{}v1/\allowbreak{}waitlist} grava o contato do interessado no PostgreSQL; reenvios com o mesmo WhatsApp retornam conflito (HTTP 409). Durante esta rodada de validação, foram removidos dois fluxos de autoatendimento que haviam sido implementados de forma parcial e que não refletiam o modelo de operação real do produto: uma tela de \textit{onboarding} público no painel (\texttt{/\allowbreak{}auth/\allowbreak{}onboarding}), que permitia a qualquer visitante criar uma nova empresa sem intermediação, e um \textit{modal} de \textit{checkout} na página institucional, integrado a um fluxo de pagamento via Mercado Pago que nunca chegou a ser conectado à interface (os botões de assinatura da seção de planos direcionavam para a lista de espera, e não para o \textit{checkout}). Em ambos os casos, o \textit{endpoint} \texttt{/\allowbreak{}api/\allowbreak{}v1/\allowbreak{}auth/\allowbreak{}provision}, responsável por criar \textit{tenant}, administrador e configurações iniciais em uma única operação, permanece a forma correta de provisionar uma nova empresa, porém acionado diretamente pelo operador da plataforma, e não por autoatendimento público. Essa decisão e sua eventual automatização futura são discutidas na Seção~\ref{sec:limitacoes} do Capítulo~\ref{cap5Discussao}.

\subsection{Módulo de Ideário (Sugestões e Roteiro de Funcionalidades)}

Durante a validação, identificou-se que o módulo de ideário estava parcialmente implementado. Esse módulo permite à equipe registrar itens de roteiro (\textit{backlog}, em desenvolvimento, lançado) e permite a qualquer usuário autenticado enviar sugestões de melhoria. No momento da validação, o domínio, o repositório, o caso de uso e o \textit{handler} HTTP já existiam no \textit{backend}, e o serviço de consumo já existia no \textit{frontend}, mas as rotas nunca haviam sido registradas no roteador da API, e a página correspondente do painel (\texttt{/\allowbreak{}dashboard/\allowbreak{}ideario}) nunca havia sido criada, apesar do item já constar no menu lateral. O resultado observado era um \textit{link} que retornava HTTP 404.

A funcionalidade foi concluída com o registro das seguintes rotas:

\begin{itemize}
  \item \texttt{GET} e \texttt{POST} \texttt{/\allowbreak{}api/\allowbreak{}v1/\allowbreak{}roadmap};
  \item \texttt{PUT} e \texttt{DELETE} \texttt{/\allowbreak{}api/\allowbreak{}v1/\allowbreak{}roadmap/\allowbreak{}:id};
  \item \texttt{POST} \texttt{/\allowbreak{}api/\allowbreak{}v1/\allowbreak{}roadmap/\allowbreak{}suggestions}.
\end{itemize}

A listagem ficou aberta a qualquer usuário autenticado, as mutações restritas ao papel \textit{admin} e envio de sugestões aberto a qualquer usuário do \textit{tenant}. A página do painel foi implementada com um quadro por status e formulários de criação, edição e sugestão. Após a conclusão, o fluxo completo foi validado com um item real já existente no banco de dados (semeado em uma etapa anterior do projeto e, até então, inacessível pela ausência da rota).

\section{Ambiente de Testes e Dispositivos Utilizados}

Complementarmente aos testes funcionais descritos na Seção~\ref{sec:testes-funcionais}, o painel foi acessado em navegadores \textit{desktop} e em dispositivos móveis físicos, na mesma rede local do ambiente de desenvolvimento, a fim de observar o comportamento real de \textit{layout}, toque e navegação.

\subsection{Dispositivos Utilizados}

A Tabela~\ref{tab:dispositivos-testados} relaciona os dispositivos e navegadores utilizados nesta rodada de testes.

% TODO(Marcelo): completar a coluna "Dispositivo / SO" das linhas iOS e Android
% com o modelo exato e a versão do sistema, por exemplo: "iPhone 13 (iOS 18)" e
% "Samsung Galaxy S23 (Android 14)".
\begin{table}[H]
  \centering
  \caption{Dispositivos e navegadores utilizados nos testes manuais de responsividade.}
  \label{tab:dispositivos-testados}
  \begin{tabular}{p{3cm}p{5.5cm}p{5.5cm}}
    \hline
    \textbf{Plataforma} & \textbf{Dispositivo / SO} & \textbf{Navegador} \\
    \hline
    Desktop & macOS, resoluções widescreen ($>$1400px) e notebook ($\sim$1280px) & Chrome e Safari \\
    iOS & [\textit{completar com modelo do iPhone e versão do iOS}] & Safari \\
    Android & [\textit{completar com modelo do aparelho e versão do Android}] & Chrome \\
    \hline
  \end{tabular}
  \begin{minipage}{0.95\textwidth}
    \vspace{0.3em}
    {\small Fonte: O Autor, 2026.}
  \end{minipage}
\end{table}

\section{Responsividade e Acessibilidade}

\subsection{Adaptação a Diferentes Resoluções}

O painel utiliza \textit{layout} responsivo baseado em Tailwind CSS, reorganizando cartões de indicadores, tabelas e formulários em coluna única em telas estreitas. As Figuras~\ref{fig:mobile-dashboard} e~\ref{fig:mobile-clientes} (Capítulo~\ref{cap3MatMet}) ilustram esse comportamento no \textit{dashboard} e no módulo de clientes.

\subsection{Tamanho de Fonte, Contraste e Área de Toque}

% TODO(Marcelo): registrar observações concretas dos testes em dispositivo real -- por exemplo,
% se algum texto ficou pequeno demais, se algum botão/ícone ficou difícil de tocar, se as tabelas
% precisaram de rolagem horizontal, etc. Um achado real aqui fortalece o capítulo.
Durante a navegação em dispositivo móvel real, observou-se que [\textit{completar com as observações da rodada de testes em iPhone/Android}]. Tabelas com muitas colunas (como as dos módulos de Clientes e Orçamentos) mantêm rolagem horizontal em telas estreitas, uma vez que não foi implementada, nesta versão, uma visualização alternativa em cartões para esses módulos em telas pequenas; esse ponto está registrado como limitação na Seção~\ref{sec:limitacoes} do Capítulo~\ref{cap5Discussao}.

\subsection{Mensagens de Acessibilidade e Leitores de Tela}

Não foi conduzida, nesta etapa, uma bateria dedicada de testes com leitores de tela ou navegação por teclado; essa lacuna é registrada como limitação na Seção~\ref{sec:limitacoes} do Capítulo~\ref{cap5Discussao} e como oportunidade de melhoria futura.

\section{Tratamento de Erros e Mensagens de Falha}
\label{sec:tratamento-de-erros}

Esta seção documenta três falhas reais identificadas durante a validação em ambiente local, com as respectivas causas-raiz e correções (ou, quando ainda pendente de correção, o diagnóstico obtido). Para cada falha, descreve-se o que foi observado, por que ela ocorreu e como foi corrigida. Optou-se por registrá-las de forma explícita porque mostram que parte relevante dos defeitos de um sistema só aparece quando ele é executado de ponta a ponta, e não na leitura isolada do código, o que justifica a etapa de validação descrita neste capítulo.

\subsection{Painel inacessível por importação órfã}
\label{sec:falha-importacao-orfa}

Ao subir o painel localmente para validação, a aplicação não renderizava nenhuma rota autenticada, falhando com erro de módulo não encontrado (\texttt{Module not found: Can't resolve '@/\allowbreak{}components/\allowbreak{}billing/\allowbreak{}SubscriptionGate'}). A investigação do histórico do repositório mostrou que um \textit{commit} anterior havia removido, de forma deliberada, toda a camada de controle de assinatura (\textit{billing}), decisão coerente com o escopo acadêmico do projeto, em que o acesso permanece aberto. Esse \textit{commit}, porém, havia deixado um \textit{import} e um uso órfãos do componente removido no \textit{layout} do painel. A correção consistiu em remover a referência remanescente, restabelecendo o funcionamento de toda a área autenticada.

\textbf{Por que ocorreu.} A remoção foi feita apagando os arquivos do componente sem procurar todos os pontos do código que ainda o utilizavam, e o erro passou despercebido porque o \textit{pipeline} do painel não executava a compilação completa (\textit{build}) nem a checagem de tipos do TypeScript antes de aceitar a alteração. No Next.js em modo de desenvolvimento, cada página só é compilada quando é acessada; como o \textit{layout} com o \textit{import} órfão só é carregado após o \textit{login}, a falha apareceu apenas quando a área autenticada foi aberta. Uma etapa de \textit{build} automática no \textit{pipeline} teria detectado o problema no momento do \textit{commit}, ponto retomado nas limitações (Seção~\ref{sec:limitacoes}).

\subsection{Cadastro de cliente sem efeito (corrigido)}
\label{sec:falha-cadastro-cliente}

Durante uma nova rodada de validação em ambiente local, observou-se que o formulário de novo cliente do painel falhava silenciosamente ao salvar: a requisição \texttt{POST /\allowbreak{}api/\allowbreak{}v1/\allowbreak{}clients} retornava HTTP 400, com o corpo da resposta indicando erros de validação para os campos \texttt{DocumentType} e \texttt{TenantID}, ambos marcados como obrigatórios (\textit{tag} \texttt{required}) no \textit{DTO} de criação do \textit{backend}. A inspeção do corpo da requisição (\textit{payload}) enviado pelo \textit{frontend} confirmou que o formulário do painel não incluía esses dois campos: o tipo de documento (CPF ou CNPJ) não era inferido nem solicitado na interface, e o identificador do \textit{tenant} não era anexado ao corpo da requisição, ainda que já estivesse disponível no \textit{token} de autenticação. O \textit{modal} apenas exibia uma mensagem genérica de falha, sem indicar a causa ao usuário. A criação de um registro de teste só foi possível, nesse momento, chamando a API diretamente com os dois campos preenchidos manualmente. A correção foi então aplicada: o formulário passou a exigir e validar explicitamente o campo \textit{Tipo} (CPF/CNPJ), com verificação de dígito verificador no próprio \textit{frontend} antes do envio (\texttt{lib/\allowbreak{}validation/\allowbreak{}clientSchema.ts}), e o \texttt{tenant\_id} passou a ser anexado automaticamente a toda requisição de criação/atualização por uma função utilitária (\texttt{withTenantId}, em \texttt{lib/\allowbreak{}api/\allowbreak{}services.ts}) que o lê do contexto de autenticação, persistido no navegador após o \textit{login}. A correção foi revisada estaticamente contra o \textit{DTO} do \textit{backend} (\texttt{CreateClientDTO}, que exige exatamente os mesmos dois campos) e contra a suíte de \textit{lint} do \textit{frontend} (\texttt{eslint}), ambos sem apontamentos; a repetição do teste manual de ponta a ponta no painel, para confirmação final, fica registrada como verificação a ser reexecutada pelo autor antes da apresentação.

\textbf{Por que ocorreu.} O \textit{backend} e o \textit{frontend} evoluíram em repositórios separados e não compartilham uma definição única do formato das requisições. Quando os campos \texttt{DocumentType} e \texttt{TenantID} passaram a ser obrigatórios no \textit{backend}, o formulário do painel não foi atualizado, e nenhum teste automatizado exercitava a chamada real entre as duas aplicações, de modo que a divergência só apareceu no uso. Além disso, exigir o \texttt{tenant\_id} no corpo da requisição é redundante, pois a API já o obtém do \textit{token} (Seção~\ref{sec:controle-acesso}) e rejeita valores diferentes; uma melhoria futura é a API preencher esse campo sozinha, eliminando a possibilidade desse tipo de erro. O problema foi agravado pela mensagem genérica do \textit{modal}, que escondia do usuário a causa da falha.

\subsection{Risco de infraestrutura mapeado e armazenamento local como mitigação}
\label{sec:falha-infra-armazenamento}

Na configuração inicial do projeto, a ausência da variável \texttt{GCS\_BUCKET} e de credenciais de conta de serviço válidas fazia com que o \textit{endpoint} de assinatura de \textit{upload} (\texttt{POST} \texttt{/\allowbreak{}api/\allowbreak{}v1/\allowbreak{}uploads/\allowbreak{}sign}) respondesse HTTP 503 em ambiente local, comportamento esperado e coberto por tratamento de erro no \textit{frontend}, mas que deixava o \textit{upload} de mídias indisponível sem uma conta de nuvem ativa. Durante a revisão do fluxo de \textit{upload} assinado e do roteiro de \textit{bootstrap} da máquina virtual (Seção~\ref{sec:workflow-deploy}), identificou-se também um risco de infraestrutura relevante para quando a operação em nuvem for retomada: a assinatura de URLs (\textit{signed URLs}) no Google Cloud Storage depende da API \textit{IAM Service Account Credentials} estar habilitada no projeto, do papel \texttt{iam.serviceAccountTokenCreator} estar concedido à conta de serviço da instância, e do escopo OAuth da instância incluir \texttt{cloud-platform} (o escopo padrão de leitura, \texttt{devstorage.read\_only}, não é suficiente). Esse mapeamento foi registrado no roteiro de \textit{bootstrap} da VM como pré-requisito explícito, de modo a evitar que a mesma causa-raiz, já conhecida a partir de experiências anteriores do autor com a mesma pilha de infraestrutura, volte a se manifestar quando a implantação em nuvem for retomada.

Como resposta direta a essa dependência, foi implementado um \textit{backend} de armazenamento alternativo em disco local (\texttt{internal/\allowbreak{}infra/\allowbreak{}storage/\allowbreak{}local}, Capítulo~\ref{cap3MatMet}), selecionado automaticamente quando \texttt{GCS\_BUCKET} não está configurado. O \textit{endpoint} de assinatura deixou de depender de qualquer credencial de nuvem: os \textit{tokens} de \textit{upload} passaram a ser assinados localmente com HMAC-SHA256 e o arquivo é gravado diretamente pela própria API. Com isso, o \textit{upload} de logo, avatar e imagens de produto passou a funcionar de ponta a ponta em ambiente local, sem custo e sem a exposição HTTP 503 anteriormente observada; o \textit{backend} do GCS permanece implementado e disponível, sendo ativado automaticamente ao configurar \texttt{GCS\_BUCKET}, caso a operação em nuvem seja retomada.

\subsection{Discussão do processo de depuração}

Em nenhum dos três casos o erro dependia de dados de entrada incomuns: os três se manifestavam de forma determinística, a cada tentativa. Os dois primeiros só se tornaram visíveis quando o fluxo correspondente foi exercitado de ponta a ponta (subida completa do painel; preenchimento e envio real do formulário de cliente), e não por inspeção isolada de código ou por testes unitários; o terceiro foi identificado por revisão do roteiro de infraestrutura, antes de qualquer tentativa de implantação em nuvem. Esse resultado reforça a importância de validar não apenas unidades de código isoladamente, mas o comportamento integrado do sistema, incluindo a configuração de infraestrutura subjacente (permissões de nuvem, variáveis de ambiente, escopos de autenticação de serviço), ponto retomado na discussão do Capítulo~\ref{cap5Discussao}. Dos três casos, dois foram corrigidos diretamente nesta etapa do trabalho: o cadastro de cliente (Seção~\ref{sec:falha-cadastro-cliente}) e a dependência de credenciais de nuvem para o \textit{upload} de mídias, substituída por um \textit{backend} de armazenamento local (Seção~\ref{sec:falha-infra-armazenamento}). O terceiro, o risco de permissões de IAM para a assinatura de URLs no Google Cloud Storage, permanece documentado como pré-requisito para quando a operação em nuvem for retomada, por não ser aplicável ao ambiente local adotado nesta fase.

\section{Desempenho da Aplicação}
\label{sec:desempenho}

O desempenho da API foi avaliado com a ferramenta de testes de carga k6, utilizando o roteiro \texttt{k6/\allowbreak{}erp-smoke.js} já existente no repositório do painel, contra os \textit{endpoints} de listagem de clientes, produtos e orçamentos, autenticados via JWT. O teste foi executado com dez usuários virtuais simultâneos, por vinte segundos, contra o ambiente local (API e banco no mesmo host, sem latência de rede pública).

\begin{itemize}
  \item Requisições totais: 2\,782, com 927 iterações completas do roteiro;
  \item Taxa de falha: 0,00\% (nenhuma requisição falhou);
  \item Tempo de resposta médio: 12,55~ms; mediana: 11,61~ms;
  \item Percentil 90: 22,14~ms; percentil 95: 26,23~ms; percentil 99: 48,99~ms;
  \item Vazão: aproximadamente 138 requisições por segundo.
\end{itemize}

Os limites definidos no próprio roteiro de teste (percentil 95 abaixo de 800~ms e percentil 99 abaixo de 1\,500~ms) foram amplamente atendidos, com margem superior a 20 vezes no percentil 95. Por se tratar de um ambiente local, sem latência de rede pública nem concorrência de outros processos na mesma instância, esses números representam um piso otimista: o comportamento sob condições reais de rede e sob a capacidade computacional de uma instância de baixo custo (como a \texttt{e2-micro} considerada na arquitetura do Capítulo~\ref{cap3MatMet}) não foi medido nesta fase, e é registrado como limitação na Seção~\ref{sec:limitacoes} do Capítulo~\ref{cap5Discussao}.

\section{Qualidade do Código e Componentização}
\label{sec:qualidade-codigo}

Como indicador objetivo de qualidade, foram executadas ferramentas de análise estática nos dois repositórios principais: \texttt{golangci-lint} no \textit{backend} (agregando os analisadores \texttt{errcheck}, \texttt{ineffassign}, \texttt{staticcheck} e \texttt{unused}, entre outros) e \texttt{eslint} no \textit{frontend} (configuração \texttt{eslint-config-next}, com as regras de \textit{hooks} do React e \texttt{@typescript-eslint}).

No estado em que o projeto se encontrava no início desta rodada de validação, o \textit{backend} apresentava 19 ocorrências (8 retornos de erro não verificados, 3 atribuições inefetivas em testes, 5 sugestões de simplificação e 3 funções não utilizadas) e o \textit{frontend} apresentava 18 ocorrências (12 erros, majoritariamente uso do tipo \texttt{any}, e 6 avisos, incluindo uma chamada de atualização de estado de forma síncrona dentro de um efeito React). Todas as ocorrências foram corrigidas: os retornos de erro passaram a ser tratados ou descartados explicitamente, os tipos \texttt{any} foram substituídos por tipos concretos ou por \texttt{Record\textless string, unknown\textgreater}, o código morto foi removido, e os dois efeitos que atualizavam estado de forma síncrona foram reestruturados: um com \texttt{useSyncExternalStore}, para a detecção de montagem do componente sem divergência de hidratação (\textit{hydration mismatch}), e outro como uma função assíncrona nomeada invocada a partir do efeito. Ao final, ambos os repositórios passaram a apresentar zero ocorrências nas duas ferramentas, com a suíte de testes automatizados do \textit{backend} (\texttt{go test ./\allowbreak{}...}) permanecendo integralmente aprovada após as correções.

\section{Integração Contínua e Preparação para Entrega Contínua}
\label{sec:ci-cd-resultados}

O \textit{pipeline} de integração contínua (GitHub Actions) foi revisado e ampliado durante esta etapa. Na configuração original, o \textit{workflow} de \textit{deploy} construía a imagem Docker, publicava-a no GHCR com duas \textit{tags} (\texttt{latest} e o \textit{hash} do \textit{commit}) e a implantava em uma máquina virtual via SSH, sem verificação posterior de sucesso. O \textit{workflow} foi ampliado para: (i) sincronizar automaticamente os arquivos de configuração do Docker Compose e do Nginx a partir do repositório, eliminando a necessidade de cópia manual; (ii) executar um \textit{health check} contra o \textit{endpoint} \texttt{/\allowbreak{}health} após a subida dos novos contêineres; e (iii) reverter automaticamente para a \textit{tag} anteriormente implantada, registrada em um arquivo de controle, caso o \textit{health check} falhe, encerrando a execução com status de falha para sinalização da equipe. Esse fluxo completo está implementado e versionado no repositório (Seção~\ref{sec:workflow-deploy}).

A etapa de construção e publicação da imagem no GHCR foi validada com execuções reais, disparadas por \textit{push} na ramificação principal, confirmando a geração correta da imagem versionada pelo \textit{hash} do \textit{commit}. A etapa final de implantação (\textit{pull} da imagem, atualização dos contêineres via SSH e \textit{health check} com reversão automática), por sua vez, foi validada isoladamente em ambiente local, sem alvo de implantação público, conforme discutido na Seção~\ref{sec:limitacoes} do Capítulo~\ref{cap5Discussao}: sua execução contínua contra uma instância de nuvem real permanece como próximo passo natural do projeto, e não como resultado obtido nesta fase acadêmica.


\section{Síntese dos Testes Executados}
\label{sec:checklist-testes}

A Tabela~\ref{tab:checklist-testes} reúne, em um único lugar, todos os cenários testados neste capítulo, para que o leitor possa conferir rapidamente o que foi verificado e qual foi o resultado de cada verificação. Cada linha corresponde a um cenário descrito nas seções anteriores, indicado na última coluna. A coluna ``Resultado'' usa três valores: \textbf{Aprovado}, quando o cenário funcionou como esperado; \textbf{Corrigido}, quando o cenário falhou na primeira execução, a causa foi identificada e corrigida, e o cenário passou a funcionar; e \textbf{Não realizado}, quando o teste não foi executado nesta etapa do trabalho e permanece como limitação.

\begin{table}[H]
  \centering
  \caption{Síntese dos cenários de teste executados.}
  \label{tab:checklist-testes}
  \small
  \begin{tabular}{p{7.9cm}p{2.8cm}p{2.9cm}}
    \hline
    \textbf{Cenário testado} & \textbf{Resultado} & \textbf{Onde é descrito} \\
    \hline
    \textit{Login} com e-mail e senha, emissão e renovação do JWT & Aprovado & Seção~\ref{sec:testes-funcionais} \\
    Criação de nova empresa (\textit{provision}) via API & Aprovado & Seção~\ref{sec:testes-funcionais} \\
    Isolamento de dados entre empresas (\textit{tenants}) distintas & Aprovado & Seção~\ref{sec:testes-funcionais} \\
    Cadastro, listagem, edição e exclusão de clientes & Aprovado & Seção~\ref{sec:testes-funcionais} \\
    Cadastro de produtos, incluindo envio de imagem & Aprovado & Seção~\ref{sec:testes-funcionais} \\
    Orçamentos, com atualização de status comercial e de obra & Aprovado & Seção~\ref{sec:testes-funcionais} \\
    Painel inicial (\textit{dashboard}) com indicadores e gráficos & Aprovado & Seção~\ref{sec:testes-funcionais} \\
    Lista de espera da página institucional, com recusa de cadastro repetido (HTTP 409) & Aprovado & Seção~\ref{sec:testes-funcionais} \\
    Módulo de ideário (itens do roteiro e sugestões) & Aprovado & Seção~\ref{sec:testes-funcionais} \\
    Envio de arquivo por URL assinada, com armazenamento local & Aprovado & Seção~\ref{sec:falha-infra-armazenamento} \\
    \textit{Pipeline} de CI/CD: \textit{build} e publicação no GHCR & Aprovado & Seção~\ref{sec:ci-cd-resultados} \\
    Cadastro de cliente pelo formulário do painel & Corrigido & Seção~\ref{sec:falha-cadastro-cliente} \\
    Testes automatizados do \textit{backend} (\texttt{go test}) & Aprovado & Seção~\ref{sec:qualidade-codigo} \\
    Teste de carga (k6) nos \textit{endpoints} de listagem & Aprovado & Seção~\ref{sec:desempenho} \\
    Testes de acessibilidade (leitor de tela, teclado, contraste) & Não realizado & Seção~\ref{sec:limitacoes} \\
    Implantação e teste de carga em instância de nuvem real & Não realizado & Seção~\ref{sec:limitacoes} \\
    \hline
  \end{tabular}
  \begin{minipage}{0.95\textwidth}
    \vspace{0.3em}
    {\small Fonte: O Autor, 2026.}
  \end{minipage}
\end{table}


\subsection*{Considerações Finais do Capítulo}

Este capítulo apresentou os resultados obtidos na validação do Onmarmoraria: testes funcionais dos módulos de negócio, testes de responsividade em dispositivos reais, três falhas documentadas com suas causas-raiz, métricas de desempenho sob carga, indicadores de qualidade de código e a verificação da etapa de integração contínua. O Capítulo~\ref{cap5Discussao} discute esses resultados frente aos objetivos definidos na Introdução, situando também as limitações identificadas e as melhorias futuras propostas.
